\documentclass[12pt]{article}
\usepackage[utf8]{inputenc}
\usepackage[T1]{fontenc}
\usepackage{hyperref}
\title{Santa Mar\'{\i}a Quiegolani Zapotec Language Resources for K-12 Curriculum Development}
\author{Compiled by Claude Sonnet 5.5 (claude-sonnet-5-5) for Andr\'as Kornai}
\date{}
\begin{document}
\maketitle

\section{Introduction}
Santa Mar\'{\i}a Quiegolani Zapotec (ISO 639-3 zpi; Glottocode sant1451) is an endangered Otomanguean language of Oaxaca, Mexico, with roughly 2,000 speakers. The search found linguistic descriptions, a New Testament and some archival recordings. It found no school materials and no dictionary.

\section{Primary Educational Resources}
No primers, textbooks or school programmes were found in the search results. Susan Regnier compiled a Zapotec alphabet guide (2012), according to a search summary, but the guide itself was not seen and its contents are unverified.

\section{Linguistic and Reference Materials}
\begin{itemize}
\item Wikipedia article on Quiegolani Zapotec: \url{https://en.wikipedia.org/wiki/Quiegolani_Zapotec}
\item ASJP basic wordlist (from a SIL source): \url{https://asjp.clld.org/languages/SANTA_MARIA_QUIEGOLANI_ZAPOTEC}
\item Glottolog entry for Regnier 1993, ``Quiegolani Zapotec Phonology'' (SIL Work Papers): \url{https://glottolog.org/resource/reference/id/90198}
\item Wikitongues language page: \url{https://wikitongues.org/languages/zpi/}
\end{itemize}
Cheryl Black's \emph{Quiegolani Zapotec Syntax: A Principles and Parameters Account} (SIL International, 2000) was mentioned in a search summary. No link was seen.

\section{Religious and Community Resources}
\begin{itemize}
\item Santa Mar\'{\i}a Quiegolani Zapotec New Testament (Digital Bible Society): \url{https://dbs.org/bibles/ZPIWYI}
\item Scripture Earth page: \url{https://www.scriptureearth.org/00eng.php?iso=zpi}
\end{itemize}

\section{Digital and Audio Resources}
AILLA (University of Texas) holds Zapotec survey recordings and elicitation material for Santa Mar\'{\i}a Quiegolani. They can be found by searching the AILLA site for ``Zapotec Survey - Santa Mar\'{\i}a Quiegolani''.

\section{Curriculum Development Context}
The scripture, the phonology and the syntax description give a base for orthography and grammar work. No dictionary or classroom materials were found, so a curriculum would need to build those with the community.
\end{document}